\section{定位：可被拒绝的纸面研究协议，不构成交易系统}

本章将系统化趋势的公开经济逻辑改写成一个可审计的\textbf{境内期货纸面研究协议}。协议记录某日应使用的实际合约、数据、时间戳和风险约束，并在条件不足时拒绝生成信号；研究输出不含方向预测或任何外部委托。当前每个品种的纸面目标手数均为零，\qzero{}。

\begin{sourcequalitybox}[CTA 的最低事实门槛]
连续主力价格、财经软件图表或单日行情截图均不足以建立可复核的 CTA 研究。每个研究日必须能对应到实际合约、结算价、成交量、持仓量、到期日、合约乘数、交易状态、涨跌停、保证金、限仓和数据来源。缺失其中任一关键字段时，记录 \texttt{DATA\_REJECT}；不得补充估计值。
\end{sourcequalitybox}

\section{第一道门：合约事实先于信号}

\begin{exhibitbox}[Exhibit 5：每个实际合约的纸面研究准入卡]
\small
\begin{tabularx}{\exhibitboxwidth}{L{2.45cm}X L{3.85cm}}
\toprule
\textbf{核验组} & \textbf{必须保存的事实} & \textbf{任一缺失时的处置} \\
\midrule
身份与时间 & 交易所、品种、合约代码、交易日、时区、合约到期/交割窗口、数据可得时间。 & 不使用连续图替代；该日无合约研究结果。 \\
价格与活跃度 & 官方或有权数据源的结算价、成交量、持仓量、交易状态与异常标志。 & 不推断可成交价格，不把盘中最后价当作结算价。 \\
规格与限制 & 合约乘数、最小变动、涨跌停、初始/维持保证金、限仓、交割月规定。 & 不计算手数、保证金或历史 P\&L。 \\
换月记录 & 旧合约的纸面平仓价、新合约的纸面开仓价、换月日期、规则和原因。 & 不把拼接收益当作可实现收益；停止该品种序列。 \\
账户/权利 & 资料使用权、品种准入、交易所/期货公司规则的当期核验记录。 & 不把“理论可交易”写成个人资格；保持 \qzero{}。 \\
\bottomrule
\end{tabularx}
\sourcenote{S11--S12、S15；本报告的合约数据治理协议。具体适当性、品种和账户权限须逐项另行核验。}
\end{exhibitbox}

\section{研究参数：须在看结果前冻结}

为避免“看见结果后调参数”，以下是本备忘录记录的\textbf{协议参数}。它们仅供审计，不用于形成建议的风险目标、保证金比例或实盘规则；在完整 R0--R4 资料、独立审查和后续许可尚未齐备前，所有参数的唯一当前输出仍是零手数。保留数值的理由是让未来研究者可以检验是否发生事后优化。

对每一实际合约 \(i\)，设 \(P_{i,t}\) 为该交易日已核验的结算价，\(\sigma_{i,60,t}\) 为截至 \(t\) 的 60 个有效交易日对数收益波动率。多速度趋势分数定义为：

\[
S_{i,t}=0.45\frac{\log(P_{i,t}/P_{i,t-21})}{\sigma_{i,60,t}\sqrt{21}}
+0.35\frac{\log(P_{i,t}/P_{i,t-63})}{\sigma_{i,60,t}\sqrt{63}}
+0.20\frac{\log(P_{i,t}/P_{i,t-126})}{\sigma_{i,60,t}\sqrt{126}}.
\]

\begin{dashboardbox}[冻结参数与解释]
\begin{tabularx}{\textwidth}{L{3.3cm}L{2.8cm}X}
\toprule
\textbf{字段} & \textbf{协议记录值} & \textbf{解释与禁止事项} \\
\midrule
观察速度 & 21 / 63 / 126 个有效交易日 & 三种速度分别占 \(0.45/0.35/0.20\)。不得因某一品种近期表现而临时增删窗口。 \\
波动率窗口 & 60 个有效交易日 & 仅用有权、可复核的实际合约结算序列；数据断裂不得静默填补。 \\
进入缓冲 & \(\vert S_{i,t}\vert \geq 0.35\) & 仅是纸面候选资格，非下单信号；需要合约、流动性和风险门禁同时通过。 \\
退出缓冲 & 已有纸面状态在 \(\vert S_{i,t}\vert \leq 0.15\) 时退出 & 用于减少阈值附近反复切换；不假设任何价格必然成交。 \\
时间纪律 & 收盘后生成；最早下一有效交易日才可观察处理 & 不允许使用未来结算价、盘后未知字段或同日完美成交假设。 \\
研究风险封顶 & 年化波动率参数 \(10\%\)；共同 NAV 初始保证金参数上限 \(25\%\) & 仅作测试围栏；个人配置建议、资金承诺和杠杆许可均未形成。 \\
\bottomrule
\end{tabularx}
\end{dashboardbox}

\section{手数计算必须先承认当前为零}

若未来获得完整且经审查的合约资料，纸面模型可以先计算每手的风险金额和保证金占用，再施加品种组、相关性、参与率和限仓约束。为了避免公式被误读成当前委托，本报告将当前状态明确写为：

\[
n^{\mathrm{paper}}_{i,t}=0,\qquad
\mathrm{Margin}^{\mathrm{paper}}_{i,t}=0,\qquad
\mathrm{ExternalOrder}_{i,t}=\varnothing.
\]

只有在未来独立审核确认了历史数据权利、实际合约映射、成本模型、点时规则、风险预算和纸面对账之后，研究者才可讨论一个\emph{假设的}非零纸面手数。即使届时讨论，下列约束仍须作为并列的否决条件。波动率公式不能单独构成充分条件：

\begin{enumerate}
  \item 每个品种的实际合约必须在非交割/限仓/停市风险窗口之外，且由当日规则确认；
  \item 单一品种和同一经济品种组的压力损失必须在预先冻结的共同 NAV 损失预算内；
  \item 计划成交量不得假设占用超过预设、并经数据验证的市场参与比例；
  \item 涨跌停、流动性中断、保证金上调或无法换月时，模型必须保留无法成交/被动持有的结果；
  \item 任一数据、权限、风险或合规字段失效，目标自动回零，不能用主观“例外”覆盖。
\end{enumerate}

\section{换月、空仓与失败日必须进入账本}

CTA 的可信度要求对不方便事实作完整记录；漂亮的连续曲线不足以建立可信度。换月至少涉及旧合约的处置、新合约的进入、价差、流动性、保证金和交割规则，不能作为价格序列中的无痕拼接点处理。空仓是当信号、合约事实或风险门禁不满足时的正确状态，不构成模型失败。下表规定纸面账本必须如何处理失败日。

\begin{exhibitbox}[Exhibit 6：CTA 账本的异常状态]
\small
\begin{tabularx}{\exhibitboxwidth}{L{3.1cm}X L{3.8cm}}
\toprule
\textbf{状态} & \textbf{必须记录的事实} & \textbf{允许的研究结论} \\
\midrule
数据断裂 & 缺失字段、来源、抓取时间、受影响合约和回补路径。 & \texttt{DATA\_REJECT}；不插值、不回看后补。 \\
换月 / 交割窗口 & 双腿合约、价格、日期、交割月限制、成交/持仓及规则来源。 & \texttt{ROLL\_REVIEW}；未闭合则该品种空仓。 \\
涨跌停 / 停市 & 交易状态、价格限制、保证金/限仓变化、未成交数量。 & \texttt{NO\_FILL\_ASSUMED}；不得按理想止损价结算。 \\
保证金压力 & 初始/维持保证金、可用现金、已占用保证金、压力情景。 & \texttt{RISK\_FREEZE}；不以加仓或融资修复。 \\
规则/资格不清 & 当期规则 URL、核验日期、账户权限或资料权利的缺口。 & {\scriptsize\texttt{ELIGIBILITY\_UNRESOLVED}}；维持 \qzero{}。 \\
\bottomrule
\end{tabularx}
\sourcenote{S11--S12、S15；本报告的纸面运行设计。}
\end{exhibitbox}

\begin{riskbox}[CTA 最常见的伪复现]
把单一主力连续合约、一个均线交叉、零成本回测和即时止损假设称为“机构 CTA”，会遗漏合约映射、风险缩放、保证金、交割、执行和容量这些策略本体。本报告在事实缺失时保留空白或零手数；图表相似性不足以支持可复制性结论。
\end{riskbox}
